%===============================================================================
% LLMWiki Technical Specification — Volume 04: Retrieval & Evidence Gate
%===============================================================================
\chapter{Retrieval and Evidence Gate}
\label{cha:retrieval}

This volume specifies the \emph{retrieval pipeline}---the hybrid search mechanism that fuses vector, keyword, and graph signals---and the \emph{Evidence Gate}, the mandatory verification stage that ensures every retrieved unit cited in an answer is semantically relevant to the query. This is the critical differentiator from prior RAG systems: retrieval is necessary but not sufficient; evidence verification is mandatory.

\section{Evidence Gate: Core Principle}
\label{sec:retrieval:evidence-gate}

\begin{designprinciple}[Evidence Closure]
No retrieval result reaches the synthesis stage without passing multi-stage verification against calibrated thresholds. This enforces \emph{evidence closure}: every citation in an answer traces to a verified unit.
\end{designprinciple}

A single cross-encoder score answers one question --- \emph{is this unit relevant to the query} --- and cannot by itself answer three others that ``evidence-aware'' implies: whether the unit's content is still true, whether a specific generated claim is actually entailed by the unit (as opposed to merely topically adjacent to it), and whether the answer as a whole is backed by evidence for every claim it makes, not just the claims someone happened to cite. Treating $\text{score}(q,u) \geq \tau$ as the definition of ``verified evidence'' conflates these. This volume therefore separates the gate into four distinct checks, only the first two of which run before synthesis (\S\ref{sec:retrieval:proofreading}); the other two run against the generated answer (\S\ref{sec:retrieval:synthesis}) and are what actually justify calling an answer evidence-closed:

\begin{enumerate}
  \item \textbf{Relevance Gate} (\S\ref{sec:retrieval:proofreading}, Stages 1--2): is the candidate unit relevant to the query at all. This is what a cross-encoder score measures.
  \item \textbf{Structural / Source Validity Gate} (\S\ref{sec:retrieval:proofreading}, Stage 3, extended below): is the candidate reachable/consistent in the graph, and --- where the source tracks it --- not stale or access-restricted.
  \item \textbf{Claim Support Gate} (\S\ref{sec:retrieval:proofreading}, Stage 4; \S\ref{sec:retrieval:synthesis}): does a specific claim, not just the unit as a whole, entail from the cited evidence, with an explicit contradiction check, not only an entailment one.
  \item \textbf{Evidence Coverage Gate} (\S\ref{sec:retrieval:synthesis}): after synthesis, is every atomic claim in the answer --- not only the ones the model chose to cite --- mapped to supporting evidence.
\end{enumerate}

\begin{definition}[Evidence Gate]
\label{def:evidence-gate}
The (pre-synthesis) Evidence Gate is a function $\mathcal{G}: \mathcal{Q} \times \mathcal{U}^* \times \Theta \to \mathcal{U}^*$ that takes a query $q$, a candidate set $\mathcal{R} \subseteq \mathcal{U}$, and a threshold $\tau \in \Theta$, and returns only those units $u \in \mathcal{R}$ where $\text{score}(q, u) \geq \tau$, where $\text{score}$ is computed by a cross-encoder. This definition covers the Relevance Gate only (item 1 above); \S\ref{sec:retrieval:proofreading} extends it to a cascade covering items 1--3, and \S\ref{sec:retrieval:synthesis} adds the post-synthesis Claim Support and Evidence Coverage checks (items 3--4) that a single pre-synthesis threshold cannot provide.
\end{definition}

\section{Why Not Reranking?}
\label{sec:retrieval:not-rerank}

The Evidence Gate is \textbf{not} a reranker. Key distinctions:

\begin{table}[htbp]
\centering
\caption{Reranker vs. Evidence Gate}
\label{tab:retrieval:reranker-vs-gate}
\begin{adjustbox}{max width=\textwidth}
\begin{tabular}{lcc}
\toprule
\textbf{Property} & \textbf{Reranker} & \textbf{Evidence Gate} \\
\midrule
Output & Reordered top-$k$ & Filtered subset (size $\leq k$) \\
Threshold & Implicit (top-$k$) & Explicit calibrated $\tau$ \\
Goal & Improve ranking & Enforce relevance / reject noise \\
Failure mode & Bad results still returned & Empty set $\to$ ``no evidence'' \\
Citation policy & Top-$k$ cited & Only verified cited \\
Calibration & Rarely calibrated & Mandatory calibration \\
\bottomrule
\end{tabular}
\end{adjustbox}
\end{table}

This distinction is fundamental: a reranker always returns $k$ results; the Evidence Gate may return $0$ (triggering ``insufficient evidence'' response), which is the correct behavior when no unit actually supports the query.

\section{Retrieval Pipeline}
\label{sec:retrieval:pipeline}

\begin{figure}[htbp]
\centering
\begin{adjustbox}{max width=\textwidth}
\begin{tikzpicture}[
  stage/.style={draw, rounded corners, fill=llmwikiblue!5, text width=2.6cm, minimum height=1.3cm, font=\small, align=center, inner sep=3pt},
  flow/.style={->, thick, >=stealth}
]
  % Row 1
  \node[stage] (query) {Query $q$};
  \node[stage, right=0.5cm of query] (embed) {Embed Query};
  \node[stage, right=0.5cm of embed] (hybrid) {Hybrid Search (Vec + BM25 + Graph)};
  \node[stage, right=0.5cm of hybrid] (candidates) {Candidates $\mathcal{R}$};

  \draw[flow] (query) -- (embed);
  \draw[flow] (embed) -- (hybrid);
  \draw[flow] (hybrid) -- (candidates);

  % Row 2
  \node[stage, below=1.1cm of query] (gate) {Evidence Gate (Cross-Encoder)};
  \node[stage, right=0.5cm of gate] (verified) {Verified $\mathcal{E}$};
  \node[stage, right=0.5cm of verified] (synth) {Synthesize + Cite};

  \draw[flow] (gate) -- (verified);
  \draw[flow] (verified) -- (synth);

  % Wrap from end of row 1 to start of row 2
  \draw[flow] (candidates.south) -- ++(0,-0.45) -| (gate.north);

  % Calibration feedback
  \node[stage, dashed, fill=gray!5, below=0.9cm of verified] (cal) {Calibration: isotonic / Platt / temperature, FAR null};
  \draw[flow, dashed] (gate.south) |- (cal.west);
  \draw[flow, dashed] (cal.east) -| (synth.south);
\end{tikzpicture}
\end{adjustbox}
\caption{Retrieval + Evidence Gate pipeline.}
\label{fig:retrieval:pipeline}
\end{figure}

\subsection{Stage 1: Query Embedding}
\label{sec:retrieval:query-embed}

Query text $q$ is embedded using the same model as ingestion (config: \texttt{[index.hnsw.dim]}). For multi-vector queries (e.g., ColBERT-style), each token embedding is preserved.

\subsection{Stage 2: Hybrid Search}
\label{sec:retrieval:hybrid-search}

Detailed in Vol.~03 \S\ref{sec:kg:hybrid-search}. Returns candidate set $\mathcal{R}$ of size $k_{cand}$ (default 50).

\subsection{Stage 3: Evidence Gate}
\label{sec:retrieval:evidence-gate-stage3}

\paragraph{Cross-Encoder Model.}
A transformer cross-encoder (e.g., \texttt{cross-encoder/ms-marco-MiniLM-L-6-v2}, \texttt{jinaai/jina-reranker-v2-base}) scores $(q, u)$ pairs:
\begin{equation}
\text{score}(q, u) = \text{sigmoid}\big(\text{CrossEncoder}([q; \text{SEP}; u.content])\big)
\end{equation}

\paragraph{Batch Processing.}
Candidates are scored in batches (config: \texttt{[evidence\_gate.batch\_size]}).

\paragraph{Thresholding.}
Units with $\text{score} < \tau$ are discarded. $\tau$ is calibrated per model/dataset (see \S\ref{sec:retrieval:calibration}).

\begin{algorithm}[H]
\caption{Evidence Gate}
\label{alg:evidence-gate}
\begin{algorithmic}[1]
\Require Query $q$, candidates $\mathcal{R}$, threshold $\tau$, batch size $B$
\Ensure Verified units $\mathcal{E}$
\State $\mathcal{E} \gets \varnothing$
\For{$i = 1$ to $\lceil |\mathcal{R}| / B \rceil$}
  \State $batch \gets \mathcal{R}[(i-1)B : iB]$
  \State $scores \gets \textsc{CrossEncoderBatch}(q, batch)$
  \ForEach{$(u, s) \in \textsc{Zip}(batch, scores)$}
    \If{$s \geq \tau$}
      \State $\mathcal{E} \gets \mathcal{E} \cup \{(u, s)\}$
    \EndIf
  \EndFor
\EndFor
\State \Return \textsc{SortByScoreDesc}($\mathcal{E}$)
\end{algorithmic}
\end{algorithm}

\subsection{Multi-Stage Proofreading Cascade}
\label{sec:retrieval:proofreading}

Algorithm~\ref{alg:evidence-gate} performs \emph{single-step} discrimination: one score, one threshold. Its achievable error rate is bounded by the separation of the score distributions of relevant and irrelevant units under a single measurement. This is the direct analogue of equilibrium binding discrimination in molecular recognition, where selectivity cannot exceed the ratio of dissociation constants.

Biological systems face an identical constraint and defeat it. A T cell must discriminate agonist from self peptide despite very small differences in receptor binding affinity, yet achieves error rates far below the equilibrium bound. The mechanism, due to \citet{hopfield1974kinetic} and \citet{ninio1975kinetic}, is \emph{kinetic proofreading}: the recognition event is decomposed into $n$ sequential intermediate states, each of which may irreversibly discard the candidate, and each of which is driven out of equilibrium by energy dissipation. Recent work has refined this picture, showing that parallel reaction threads on a single receptor improve discrimination fidelity without lengthening the sequential chain \citep{kirby2025parallel}, and that multimolecular schemes can escape the classical activity--fidelity trade-off \citep{multimolecular2026}.

We adopt this architecture directly.

\begin{designprinciple}[Proofreading over Thresholding]
Evidence verification is a cascade of $n$ irreversible discard stages of increasing cost and decreasing throughput, not a single thresholded score. Fidelity is purchased with computation, and the exchange rate is quantifiable.
\end{designprinciple}

\begin{definition}[Proofreading Cascade]
\label{def:proofreading-cascade}
A proofreading cascade of depth $n$ is a sequence of verifiers $g_1, \dots, g_n$ with thresholds $\tau_1, \dots, \tau_n$ and unit costs $c_1 \leq \dots \leq c_n$. A candidate $u$ survives stage $j$ iff $g_j(q,u) \geq \tau_j$, and is discarded irreversibly otherwise. The gate output is the set of candidates surviving all $n$ stages.
\end{definition}

\begin{theorem}[Multiplicative Error Suppression]
\label{thm:error-suppression}
Let $\varepsilon_j$ denote the false-acceptance probability of stage $j$ on an irrelevant unit, and $\rho_j$ its true-acceptance probability on a relevant unit. If the stage verdicts are conditionally independent given relevance, the cascade attains
\begin{equation}
\varepsilon_{\text{casc}} = \prod_{j=1}^{n} \varepsilon_j, \qquad
\rho_{\text{casc}} = \prod_{j=1}^{n} \rho_j .
\end{equation}
For homogeneous stages ($\varepsilon_j = \varepsilon$, $\rho_j = \rho$) the false-acceptance rate decays as $\varepsilon^n$ while recall decays only as $\rho^n$; since $\rho \gg \varepsilon$ by construction, the discrimination ratio $\rho_{\text{casc}}/\varepsilon_{\text{casc}} = (\rho/\varepsilon)^n$ grows exponentially in depth.
\end{theorem}

The conditional independence premise is the load-bearing assumption, and it is not free: stacking $n$ copies of the same cross-encoder yields perfectly correlated verdicts and no gain whatsoever. Independence is engineered by making the stages read \emph{different evidence channels}, which is the retrieval analogue of the parallel reaction threads of \citet{kirby2025parallel}:

\begin{enumerate}
  \item \textbf{Stage 1 --- Lexical admissibility} ($c_1$ negligible): BM25 span overlap between query terms and unit content. Rejects candidates retrieved purely by embedding drift.
  \item \textbf{Stage 2 --- Cross-encoder relevance} ($c_2$ moderate): Algorithm~\ref{alg:evidence-gate} as specified above.
  \item \textbf{Stage 3 --- Structural and source-validity admissibility} ($c_3$ low): the unit must be reachable in the knowledge graph from at least one other surviving unit within $k$ hops, or be the unique survivor (rejects semantically plausible but contextually orphaned units), \emph{and}, where the source artifact carries the relevant metadata, must not be flagged stale (\texttt{ArtifactVersion.observed\_at} older than a configured freshness window) or access-restricted for the requesting principal. This stage does not adjudicate factual correctness of the unit's content --- only whether the unit's provenance metadata makes it eligible to be cited at all.
  \item \textbf{Stage 4 --- Entailment verification} ($c_4$ high): a natural language inference model must find the candidate answer span entailed, not merely topically related, by the unit, and separately scores the contradiction class so an evidently contradicting unit is discarded rather than merely scored as not-entailing.
\end{enumerate}

Stages 1--4 are candidate-side: they decide which retrieved units are eligible to be shown to the synthesis model. None of them inspect what the model actually writes. \S\ref{sec:retrieval:synthesis} adds the two checks that do.

\paragraph{Cost.} Naively a cascade is $n$ times more expensive. It is not, because survival is exponentially rare: with per-stage survival fraction $s_j$ the expected cost per initial candidate is
\begin{equation}
C = c_1 + \sum_{j=2}^{n} c_j \prod_{i<j} s_i ,
\end{equation}
so the expensive terminal stages are evaluated on a vanishing fraction of the candidate set. Ordering stages by increasing $c_j / (1 - s_j)$ minimises $C$ --- cheap, highly selective filters first.

\begin{observation}[Fidelity Requires Dissipation]
\label{obs:dissipation}
In the thermodynamic setting, proofreading fidelity beyond the equilibrium bound is possible only if the cycle consumes free energy; a cascade run at equilibrium provides no discrimination gain. The computational analogue is exact: the discard steps are irreversible, and the cascade's advantage over a single threshold is paid for strictly in compute. This yields an accuracy-versus-compute frontier for evidence verification rather than an unbounded accuracy claim, and we regard the existence of that frontier as a design constraint rather than a limitation to be hidden.
\end{observation}

\subsection{Stage 5: Evidence Context Selection and WC/1 Transport}
\label{sec:retrieval:context-shaping}

The gate output $\mathcal{E}$ is admissible but may exceed the reader budget
$B$. Context construction is therefore a constrained selection problem, not a
license to summarize or rewrite evidence. The safe baseline sorts by gate score
and takes the largest prefix that fits. The experimental reference policy
maximises weighted concept coverage per token with deterministic
\texttt{content\_id} tie-breaking and a best-singleton guard. This follows the
budgeted submodular selection family \citep{lin2010budgeted,peng2025adagres};
recent multi-hop RAG evidence shows that packing gains are conditional on
retrieval recall, complementary evidence structure, budget tightness, and
reader capacity \citep{bala2026survives}. LLMWiki therefore retains top-$k$ as
the production default until the A/B protocol in Vol.~06 shows an end-to-end
quality gain.

The selected set $\mathcal{E}_B$ crosses the model boundary through the
line-oriented \texttt{WC/1} evidence envelope. This removes repeated structural
keys and instruction ``prompt tax'' without compressing the evidence itself
\citep{matveev2026toon}:

\begin{verbatim}
WC/1 evidence envelope. A line '@<path>' names a source file...

Q How is currency validated?
@src/payment.py
0 code 3
def process_payment(amount, currency):
    validate_currency(currency)
    return gateway.charge(amount)
@src/validator.py
1 code 2
def validate_currency(currency):
    ...
\end{verbatim}

The prefix is query-independent and may be prompt-cacheable. Each block header
is \texttt{<id> <kind> <line-count>}; the parser then reads exactly that many
lines, so evidence containing \texttt{@decorator}, a header-shaped line, or
arbitrary Markdown remains unambiguous. The model sees only response-local
integer IDs. The host retains
\[
M:\texttt{id}\mapsto
(\texttt{content\_id},\texttt{occurrence\_id},\texttt{document},\texttt{span})
\]
and resolves model citations through $M$ after generation. Evidence text is
verbatim; normalization, summarization, and abbreviation are forbidden.
Score ordering remains deterministic, but the evaluation protocol must rotate
or otherwise control evidence position because long-context readers can
underuse evidence placed in the middle \citep{liu2024lostmiddle}.

\paragraph{Measured transport cost.}
Using \texttt{tiktoken o200k\_base}, the executable reference implementation
reduces the complete cold prompt by 30.3\% on the 20-block English fixture and
32.8\% on the 8-block Korean fixture against its JSON fallback. With the fixed
52-token prefix cached, the reductions are 31.7\% and 36.6\%. Earlier
35.8\%/41.8\% figures excluded the contract prefix and used a more verbose JSON
baseline; they are retained only in the design record, not as official results.
The transport result establishes token reduction only. Whether the model
understands \texttt{WC/1} as reliably as JSON is an open A/B question.

\paragraph{Boundary validation and fallback.}
Queries, source paths, and kind labels occupy single-line grammar positions and
are validated before serialization; block IDs and line counts must be
non-negative and IDs unique. Unknown contract versions, unsupported
capabilities, or validation failure select the JSON path. The contract is an
optimization, never a correctness requirement.

\subsection{Stage 6: Synthesis, Then Claim Support and Coverage Verification}
\label{sec:retrieval:synthesis}

This prompt \emph{requests} that the LLM ground every claim in the provided evidence; it does not \emph{guarantee} it. A model can still cite a source next to a claim the source does not actually support, or state an uncited claim entirely from parametric knowledge. Treating the prompt instruction as sufficient for evidence closure is exactly the gap the Relevance Gate (candidate-side) leaves on the synthesis side, so it is closed by inspecting the output, not by trusting the instruction:

\begin{enumerate}
  \item \textbf{Claim decomposition.} The generated answer $a$ is segmented into atomic claims $\{c_1, \dots, c_m\}$ (one checkable factual statement each).
  \item \textbf{Claim Support Gate.} For each claim $c_i$ and its cited evidence unit(s), the same NLI model used in Stage 4 scores $\text{entail}(c_i, u)$ and $\text{contradict}(c_i, u)$. A claim with no cited unit scoring above the entailment threshold, or with any cited unit scoring above the contradiction threshold, fails this gate.
  \item \textbf{Evidence Coverage Gate.} Coverage is $|\{c_i : c_i \text{ passes Claim Support}\}| / m$. An answer below a configured coverage threshold is not returned as-is: failing claims are either re-synthesized against $\mathcal{E}$, dropped with an explicit ``unsupported'' marker, or the whole answer is downgraded to ``insufficient evidence,'' depending on deployment policy (\S\ref{sec:retrieval:evidence-config}).
\end{enumerate}

The mapping produced by this stage --- $\text{claim\_id} \to \{\text{supporting evidence occurrence\_ids}\} \to \text{entailment score} \to \text{contradiction score}$ --- is what \texttt{ApprovalStore}/citation UI surfaces to a human reviewer; it is a stronger and independently auditable claim than ``the prompt asked for citations.''

\section{Threshold Calibration: Supervised Score Calibration}
\label{sec:retrieval:calibration}

The threshold $\tau$ must be calibrated to achieve desired precision/recall tradeoff. Uncalibrated cross-encoder scores are not probabilities. Unlike \S\ref{sec:retrieval:far}'s Negative-Only FAR Estimation, every method below requires a labeled validation set with positive examples --- this is \emph{Supervised Score Calibration} in the terminology of \S\ref{sec:retrieval:far}, named separately because it answers a different question (what is the calibrated probability of relevance) and has a different cost (annotation) than the label-free procedure.

\subsection{Calibration Methods}
\label{sec:retrieval:cal-methods}

\begin{enumerate}
  \item \textbf{Isotonic Regression}: Non-parametric, fits piecewise constant function. Requires validation set $\mathcal{D}_{val} = \{(q, u, y)\}$ where $y \in \{0,1\}$ is relevance judgement.
  \item \textbf{Platt Scaling}: Fits logistic regression $\sigma(As + B)$ to scores $s$. Simpler, less flexible.
  \item \textbf{Temperature Scaling}: Single parameter $T > 0$, calibrated $\sigma(s / T)$. Most efficient.
\end{enumerate}

\begin{algorithm}[H]
\caption{Calibration Procedure}
\label{alg:calibration}
\begin{algorithmic}[1]
\Require Validation set $\mathcal{D}_{val} = \{(q_i, u_i, y_i)\}$, cross-encoder $f$
\Ensure Calibrated threshold $\tau^*$ for target precision $P_{target}$
\State $scores \gets [f(q_i, u_i) \text{ for } (q_i, u_i, y_i) \in \mathcal{D}_{val}]$
\State $labels \gets [y_i \text{ for } (q_i, u_i, y_i) \in \mathcal{D}_{val}]$
\State Fit calibrator $C$ on $(scores, labels)$ \Comment{Isotonic/Platt/Temperature}
\State $cal\_scores \gets C(scores)$
\State $\tau^* \gets \text{threshold achieving } P_{target} \text{ on } \mathcal{D}_{val}$
\State \Return $\tau^*, C$
\end{algorithmic}
\end{algorithm}

Calibration artifacts are stored in \texttt{.llmwiki/evidence/cross-encoder/calibration.json} and loaded at startup.

\subsection{Empirical Null Thresholding: Negative-Only FAR Estimation}
\label{sec:retrieval:far}

This section, \S\ref{sec:retrieval:cal-methods} (\emph{Supervised Score Calibration}: isotonic/Platt/temperature, requires labeled $\mathcal{D}_{val}$), and \S\ref{sec:retrieval:confidence} (\emph{Coverage Calibration}: conformal prediction) are three procedures that all produce a number that looks like a calibrated confidence, and are not interchangeable. This section's procedure --- discordant-pairing against an empirical null --- is named \emph{Negative-Only FAR Estimation} throughout to keep it distinct from the other two: it is built entirely from constructed negatives (Definition~\ref{def:discordant-pair}) and answers only ``how often would a score this high occur from unrelated evidence,'' never ``what fraction of admitted units are actually relevant'' (that is precision, which needs the positive-side distribution) or ``what fraction of true positives survive'' (that is recall/FDR, which needs labeled positives too). The paragraph below states this limitation explicitly rather than letting the similarity of the numbers invite the conflation.

Algorithm~\ref{alg:calibration} requires $\mathcal{D}_{val}$ with human relevance judgements $y_i$. This is the binding practical constraint on every deployment: relevance labels are expensive, corpus-specific, and stale the moment the corpus changes. A threshold that cannot be recalibrated without re-annotation is a threshold that will silently drift.

Gravitational-wave astronomy confronts a structurally identical problem --- assigning significance to candidate events when no labelled catalogue of true detections exists --- and solves it without labels. The technique is \emph{time slides} \citep{was2010timeslides,pycbclive2018}: triggers from one detector are shifted in time by more than the inter-detector light travel time, so that any coincidence formed from the shifted streams \emph{cannot} correspond to a real astrophysical signal. Repeating the shift many times yields a large empirical sample of pure-noise coincidences, from which the false alarm rate (FAR) is read off as a function of the ranking statistic. Detection thresholds are then quoted in operational units --- LIGO's public alert threshold corresponds to approximately one false alarm per two months.

We construct the same empirical null for evidence scoring.

\begin{definition}[Discordant Pair]
\label{def:discordant-pair}
A discordant pair $(q, u^-)$ is a query paired with a unit drawn such that genuine relevance is excluded by construction, not by judgement. Admissible constructions are:
\begin{enumerate}
  \item \textbf{Corpus displacement}: $u^-$ is drawn from a repository disjoint in domain from the one $q$ was issued against.
  \item \textbf{Temporal displacement}: $u^-$ is drawn from a WAL snapshot predating the creation of every entity referenced by $q$, so no unit in that snapshot can describe them.
  \item \textbf{Graph displacement}: $u^-$ is drawn from a connected component of the knowledge graph containing no entity mentioned in $q$.
\end{enumerate}
\end{definition}

\begin{algorithm}[H]
\caption{False Alarm Rate Estimation without Relevance Labels}
\label{alg:far}
\begin{algorithmic}[1]
\Require Query sample $Q$, corpora $\{\mathcal{U}_1,\dots,\mathcal{U}_m\}$, scorer $f$, displacement count $N$
\Ensure Empirical null $\hat{F}_0$, threshold $\tau^*$ for target $\text{FAR}^*$
\State $S_0 \gets \varnothing$
\For{$i = 1$ to $N$}
  \ForAll{$q \in Q$}
    \State $u^- \gets \textsc{DrawDiscordant}(q, \{\mathcal{U}_j\})$ \Comment{Def.~\ref{def:discordant-pair}}
    \State $S_0 \gets S_0 \cup \{f(q, u^-)\}$
  \EndFor
\EndFor
\State $\hat{F}_0(s) \gets |\{s' \in S_0 : s' \geq s\}| \, / \, |S_0|$ \Comment{Empirical survival function}
\State $\tau^* \gets \min \{ s : \hat{F}_0(s) \cdot \bar{k} \leq \text{FAR}^* \}$ \Comment{$\bar{k}$ = mean candidates per query}
\State \Return $\hat{F}_0, \tau^*$
\end{algorithmic}
\end{algorithm}

\paragraph{Reporting.} Every admitted unit carries not a raw score but a false alarm rate: \emph{``evidence of this strength arises from an unrelated corpus once per $1/(\hat{F}_0(s)\bar{k})$ queries.''} This is an operationally meaningful quantity that a threshold of $0.72$ is not.

\begin{invariant}[Operational Significance]
\label{inv:far}
No unit enters the synthesis context without an associated FAR computed against the current empirical null. Recalibration is triggered whenever the WAL advances past a configured number of commits, and requires no human annotation.
\end{invariant}

\paragraph{Limitations, stated plainly.} Two are inherited directly from the source domain and one is specific to ours.
\begin{itemize}
  \item \textbf{Null contamination.} Time slides work because the light travel time bound is a physical guarantee. Our displacements are weaker guarantees: a unit from a ``disjoint'' repository may still be generically relevant (a logging idiom, a licence header). Contaminated discordant pairs inflate the null's upper tail and make the gate conservative. Domain disjointness must be asserted from repository metadata, never assumed.
  \item \textbf{Estimator saturation.} The gravitational-wave literature reports that FAR estimation error saturates with the number of time shifts \citep{was2010timeslides}; beyond a point, additional displacements re-use the same triggers and add no independent information. The same ceiling applies to $N$ in Algorithm~\ref{alg:far}, and we therefore treat $N$ as tuned to the saturation point rather than maximised.
  \item \textbf{Null $\neq$ alternative.} FAR bounds the rate of \emph{spurious} admission. It says nothing about recall. FAR calibration therefore complements, and does not replace, the labelled conformal procedure of \S\ref{sec:retrieval:confidence}: the empirical null fixes the operating point without labels, and conformal prediction supplies coverage guarantees when labels are available.
\end{itemize}

\subsection{Confidence Intervals: Coverage Calibration}
\label{sec:retrieval:confidence}

Beyond point threshold, we provide confidence intervals per prediction. This is \emph{Coverage Calibration} in the terminology of \S\ref{sec:retrieval:far}: like Supervised Score Calibration above, conformal prediction requires labeled examples, but it targets a distribution-free coverage guarantee on intervals rather than a single calibrated point score.
\begin{itemize}
  \item \textbf{Conformal Prediction}: Split conformal for distribution-free coverage guarantees
  \item \textbf{MC Dropout}: Enable dropout at inference, multiple forward passes $\to$ uncertainty
  \item \textbf{Ensemble}: Multiple cross-encoders $\to$ variance as uncertainty
\end{itemize}

Default: Temperature scaling + conformal prediction with 95\% coverage.

\section{Global Consistency of Evidence (Experimental Research Hypothesis)}
\label{sec:retrieval:sheaf}

\begin{quote}
\textit{Status: this section is a research hypothesis with a worked-out mathematical construction, not a shipped or empirically validated mechanism. It is written at specification-level precision so it is falsifiable (\S\ref{sec:retrieval:sheaf-gaps} states exactly what would falsify it and what remains unspecified), not because it has been implemented and measured. \S\ref{sec:retrieval:sheaf-gaps} and Vol.~11 \S\ref{sec:paper:results} track what validating it requires.}
\end{quote}

Everything specified so far verifies each unit \emph{independently} against the query. This leaves a hole that no amount of per-unit accuracy can close: two units may each be strongly relevant to $q$ and yet assert incompatible things about it. A gate built from pairwise scores admits both, the synthesis stage receives a contradictory evidence set, and the resulting answer is fluently grounded in mutually exclusive sources. Evidence closure as defined in \S\ref{sec:retrieval:evidence-gate} is a statement about provenance, not about coherence, and is therefore not yet closed.

The obstruction to assembling locally valid data into a globally valid whole is the subject of sheaf theory, where it is not a metaphor but the definition. A sheaf assigns data to open sets and specifies restriction maps between them; the global sections $H^0$ are exactly the consistent global assignments, and the first cohomology $H^1$ measures the obstruction to gluing local data that is pairwise compatible into a global section. Cellular sheaves make this computable over graphs, and recent work has applied the construction to distributed computation --- where terminating solutions to a task are precisely the global sections of a task sheaf \citep{sheafdistributed2025} --- to consistency constraints on networks \citep{sheafnetworks2025}, and to characterising irreducible error patterns in predictive coding networks \citep{sheafpc2025}.

\begin{definition}[Evidence Sheaf]
\label{def:evidence-sheaf}
Let $\mathcal{E}$ be the set of units surviving the proofreading cascade, and let $G_{\mathcal{E}}$ be the graph with vertices $\mathcal{E}$ and an edge $\{u,v\}$ whenever $u$ and $v$ share an entity or stand in a graph relation. Fix a coefficient field $k = \mathbb{R}$. For each vertex $u$, let $C_u$ be the finite set of atomic claims extracted from $u$ (\S\ref{sec:retrieval:synthesis}); the stalk is $\mathcal{F}(u) = k^{C_u}$, a vertex's truth-assignment \emph{simplex} embedded linearly (coordinate $c$ of a vector holds the model's confidence that claim $c$ holds, not a boolean). For each edge $\{u,v\}$, let $C_{uv} = C_u \cap_{\text{sem}} C_v$ be the claims shared between $u$ and $v$ under semantic matching (paraphrase/entity-overlap clustering, not string identity); the stalk is $\mathcal{F}(uv) = k^{C_{uv}}$. The restriction map $\mathcal{F}_{u \trianglelefteq uv} : \mathcal{F}(u) \to \mathcal{F}(uv)$ is the linear projection onto the shared coordinates, composed with the NLI model's entailment/contradiction/neutral distribution at those coordinates: entailment maps a claim's confidence forward unchanged, contradiction maps it to its negation ($x \mapsto 1-x$), and neutral maps it to $0$ (abstains the coordinate from the gluing constraint). This linearization is a modeling choice, not a derivation from the NLI model's output, and \S\ref{sec:retrieval:sheaf-gaps} states what it assumes away.
\end{definition}

\begin{definition}[Evidence Closure, Cohomological Form]
\label{def:closure-cohomological}
Let $d: C^0(G_{\mathcal{E}}; \mathcal{F}) \to C^1(G_{\mathcal{E}}; \mathcal{F})$ be the sheaf coboundary map, $H^0 = \ker d$, $H^1 = \mathrm{coker}\, d$. These are \emph{not} equivalent statements and neither alone answers the question this gate needs answered:
\begin{itemize}
  \item $H^0 \neq \varnothing$ (equivalently, $\dim \ker d > 0$) says a globally consistent vertex-assignment \emph{exists somewhere} in the stalk space --- it does not say the assignment the units actually asserted is that one.
  \item $H^1 \neq 0$ says $d$ is not surjective --- some hypothetical edge data is not induced by any global section --- which is a property of the sheaf's shape (its overlap graph and restriction maps), not of any particular observed assignment.
\end{itemize}
The question this gate actually needs is neither: given the \emph{specific} edge data $s \in C^1$ that the NLI restriction maps produced from the retrieved units, does $s$ arise from a global section? Equivalently, is $s \in \operatorname{im} d$, i.e.\ does its class $[s] \in H^1$ (via the quotient $C^1 \twoheadrightarrow \mathrm{coker}\,d$) vanish. An evidence set $\mathcal{E}$ is \emph{closed} for query $q$ iff $[s] = 0$ and the resulting global section's restriction to the answer's cited claims matches the answer; it is \emph{obstructed} iff $[s] \neq 0$, and the obstruction class's support (which coordinates of $s$ are responsible) is the conflict witness returned to the user. This is the $\delta(s) \neq 0$ test, not the unconditional claim ``$H^1(\mathcal{F}) \neq 0$''; Algorithm~\ref{alg:consistency-gate}'s \textsc{ComputeH0} step already computes this correctly (it tests the specific assignment, not sheaf-wide emptiness) and the definition above is stated to match the algorithm rather than the informal claim in earlier drafts of this volume.
\end{definition}

\begin{algorithm}[H]
\caption{Consistency Gate}
\label{alg:consistency-gate}
\begin{algorithmic}[1]
\Require Verified units $\mathcal{E}$, query $q$
\Ensure Coherent subset $\mathcal{E}^*$ or abstention with conflict witness
\State $G_{\mathcal{E}} \gets \textsc{BuildOverlapGraph}(\mathcal{E})$
\State $\mathcal{F} \gets \textsc{BuildEvidenceSheaf}(G_{\mathcal{E}})$ \Comment{Def.~\ref{def:evidence-sheaf}}
\State $Z \gets \textsc{ComputeH0}(G_{\mathcal{E}}, \mathcal{F})$ \Comment{Consistent global assignments}
\If{$Z \neq \varnothing$}
  \State \Return $\mathcal{E}$ \Comment{Closed; proceed to synthesis}
\EndIf
\State $B \gets \textsc{ComputeH1Support}(G_{\mathcal{E}}, \mathcal{F})$ \Comment{Edges carrying the obstruction}
\State $\mathcal{E}^* \gets \textsc{MaxConsistentSubset}(\mathcal{E}, B)$
\If{$\textsc{Supports}(\mathcal{E}^*, q)$}
  \State \Return $\mathcal{E}^*$ with conflict report $B$
\EndIf
\State \Return \textsc{Abstain}(``contradictory evidence'', witness $B$)
\end{algorithmic}
\end{algorithm}

\paragraph{Why cohomology rather than pairwise conflict detection.} A pairwise contradiction check finds edges where two units disagree. It cannot detect \emph{cyclic} inconsistency, in which every pair is individually compatible but no assignment satisfies the cycle --- the situation where $u$ is compatible with $v$, $v$ with $w$, and $w$ with $u$, yet the three cannot hold simultaneously. This is precisely the regime $H^1$ was invented to measure, and it is the regime in which a system reporting high pairwise confidence is most confidently wrong. The cohomology computation additionally \emph{localises} the failure: the support of the obstruction class identifies which relations carry it, so the system can report a conflict witness rather than an unexplained refusal.

\paragraph{Cost.} For the linear sheaves used here, $H^0$ and $H^1$ reduce to kernel and cokernel computations on the sheaf Laplacian, polynomial in $|\mathcal{E}|$. Since $|\mathcal{E}|$ after the cascade is small by construction (\S\ref{sec:retrieval:proofreading}), the consistency gate is not the bottleneck; the NLI evaluations that populate the restriction maps are, and they are bounded by the number of overlap edges.

\subsection{What This Construction Does Not Yet Specify}
\label{sec:retrieval:sheaf-gaps}

Definitions~\ref{def:evidence-sheaf}--\ref{def:closure-cohomological} are precise enough to implement and to falsify, which is the point of stating them formally rather than only informally. Precision makes the remaining gaps visible instead of hiding them in prose:

\begin{itemize}
  \item \textbf{Linearization is a modeling choice, not a derivation.} Truth assignment to a claim is naturally a boolean or a bounded probability, not a vector space element; embedding it in $\mathbb{R}^{C_u}$ and mapping contradiction to $x \mapsto 1-x$ is one convenient encoding among several (a simplex/statistical-manifold embedding, or working with the sheaf of sets directly via lifting to a groupoid, are alternatives with different obstruction theories). Results obtained under this encoding may not transfer to another.
  \item \textbf{Uncertainty is not propagated.} NLI models output calibrated-or-not probabilities, not certainties; the current construction takes a point estimate as the restriction map's output and does not carry the NLI model's own confidence into the obstruction computation, so a low-confidence contradiction and a high-confidence one contribute identically to $[s]$.
  \item \textbf{Only linear (additive) inconsistency is representable.} Genuine logical relationships between claims --- entailment chains, exclusive-or structure, quantified statements --- are not in general expressible as a linear restriction map; the construction here can miss a nonlinear contradiction that a symbolic or constraint-based checker would catch, and can in principle also flag a linear ``obstruction'' that a full logical reading would resolve as non-contradictory.
  \item \textbf{The overlap graph itself is a heuristic.} $G_{\mathcal{E}}$'s edges (shared entity or graph relation) determine which claims are even compared; two units that contradict each other without sharing an extracted entity produce no edge and are invisible to this gate regardless of the cohomology.
\end{itemize}

None of these gaps invalidate the construction; they are exactly what \S\ref{sec:retrieval:sheaf}'s status note asks a reader to weigh before treating cohomological closure as validated rather than specified, and what the metrics defined in Vol.~11 \S\ref{sec:paper:eval} (cyclic inconsistency detection, FAR validity) are designed to test empirically once \texttt{LLMWikiBench} exists.

\section{Query Planner Integration}
\label{sec:retrieval:planner-integration}

The Query Planner (Vol.~07) decomposes complex queries into subqueries, each executed through the retrieval + evidence gate pipeline:

\begin{verbatim}
Query: "How does the payment module handle refunds?"

Plan:
  Subquery 1: "refund processing logic" -> Search + Gate -> [E1, E2]
  Subquery 2: "refund API endpoints" -> Search + Gate -> [E3]
  Subquery 3: "refund error codes" -> Search + Gate -> [E4, E5]

Synthesis: Combine [E1..E5] with citations
\end{verbatim}

Each subquery has its own budget and can use different retrieval parameters (e.g., more graph hops for API questions, more vector weight for logic questions).

\section{Evidence Gate Config}
\label{sec:retrieval:evidence-config}

\begin{verbatim}
[evidence_gate]
model = "cross-encoder/ms-marco-MiniLM-L-6-v2"
threshold = 0.72              # calibrated threshold
calibration = "isotonic"      # "none" | "platt" | "isotonic" | "temperature"
batch_size = 16
max_candidates = 50           # max inputs to gate
max_verified = 20             # max outputs from gate
timeout_ms = 2000             # per-query gate timeout

# Conformal prediction
conformal = true
coverage = 0.95
calibration_set = "eval/evidence_calibration.jsonl"
\end{verbatim}

\section{Evaluation Metrics}
\label{sec:retrieval:metrics}

The Evidence Gate is evaluated on:

\begin{table}[htbp]
\centering
\caption{Evidence Gate Evaluation Metrics}
\label{tab:retrieval:metrics}
\begin{adjustbox}{max width=\textwidth}
\begin{tabular}{lll}
\toprule
\textbf{Metric} & \textbf{Definition} & \textbf{Target} \\
\midrule
Precision@k & $\frac{|\{u \in \text{top-}k(\mathcal{E}): y_u=1\}|}{k}$ & $> 0.85$ \\
Recall@k & $\frac{|\{u \in \mathcal{E}: y_u=1\}|}{|\{u: y_u=1\}|}$ & $> 0.75$ \\
F1 & Harmonic mean of P/R & $> 0.80$ \\
AUC-ROC & Area under ROC curve & $> 0.90$ \\
ECE (Expected Calibration Error) & $\sum_b \frac{|B_b|}{N} |acc(B_b) - conf(B_b)|$ & $< 0.05$ \\
Conformal Coverage & Empirical coverage of prediction sets & $\approx 0.95$ \\
Empty Rate & $\frac{|\{q: \mathcal{E}(q) = \varnothing\}|}{|Q|}$ & $< 0.10$ \\
\bottomrule
\end{tabular}
\end{adjustbox}
\end{table}

\section{Comparison with Prior Art}
\label{sec:retrieval:prior-art}

\begin{table}[htbp]
\centering
\caption{Retrieval + Verification in Prior Art}
\label{tab:retrieval:prior-art}
\begin{adjustbox}{max width=\textwidth}
\begin{tabular}{lccccc}
\toprule
\textbf{System} & \textbf{Hybrid} & \textbf{Cross-Enc.} & \textbf{Mandatory} & \textbf{Label-Free} & \textbf{Global} \\
 & \textbf{Search} & & \textbf{Gate} & \textbf{FAR} & \textbf{Consistency} \\
\midrule
Standard RAG & \xmark (vec only) & \xmark & \xmark & \xmark & \xmark \\
RAG with Reranker & \cmark (vec) & \cmark & \xmark & \xmark & \xmark \\
GraphRAG & \cmark (vec+graph) & \xmark & \xmark & \xmark & \xmark \\
HippoRAG & \cmark (graph PPR) & \xmark & \xmark & \xmark & \xmark \\
LightRAG & \cmark (dual-level) & \xmark & \xmark & \xmark & \xmark \\
REALM/RA-CM3 & \xmark & \xmark & \xmark & \xmark & \xmark \\
LLMWiki & \cmark (3-signal) & \cmark (cascade) & \cmark & \cmark & \cmark ($H^1$) \\
\bottomrule
\end{tabular}
\end{adjustbox}
\end{table}

The final two columns are the ones without prior occupants. No deployed retrieval system reports significance against an empirical null constructed without relevance labels, and none certifies that its returned evidence set is mutually consistent rather than merely individually relevant.

\section{Patent-Relevant Novelty}
\label{sec:retrieval:patent}

\begin{patentclaim}
\textbf{Claim 6:} A retrieval system comprising: (a) a hybrid searcher combining vector similarity, keyword matching, and graph adjacency traversal, (b) an evidence gate that applies a cross-encoder to every retrieval candidate, (c) a calibrated threshold that rejects candidates below a precision-targeted score, and (d) a synthesis stage that only cites units passing the evidence gate, ensuring evidence closure.
\end{patentclaim}

\begin{patentclaim}
\textbf{Claim 7:} The system of Claim 6 wherein the evidence gate operates as a filter (not a reranker), such that zero candidates may be returned when no evidence meets the threshold, triggering an explicit ``insufficient evidence'' response rather than hallucinated synthesis.
\end{patentclaim}

\begin{patentclaim}
\textbf{Claim 8:} The system of Claim 6 wherein the evidence gate comprises a cascade of $n \geq 2$ verification stages reading mutually distinct evidence channels, each stage irreversibly discarding candidates below a stage-specific threshold, ordered by increasing unit cost, such that the false-acceptance rate decays multiplicatively in the number of stages while the expected cost per candidate remains bounded by the survival-weighted sum of stage costs.
\end{patentclaim}

\begin{patentclaim}
\textbf{Claim 9:} The system of Claim 6 wherein the threshold is derived without relevance annotations by (a) constructing query--unit pairs for which relevance is excluded by corpus, temporal, or graph displacement, (b) accumulating scorer outputs over said pairs into an empirical null distribution, and (c) selecting the threshold attaining a target false alarm rate under said distribution, the false alarm rate being reported with each admitted unit.
\end{patentclaim}

\begin{patentclaim}
\textbf{Claim 10:} The system of Claim 6 further comprising a consistency gate that constructs a cellular sheaf over the verified evidence set with restriction maps induced by inference between overlapping claims, computes the obstruction class of the specific edge data induced by said inference (as distinct from testing whether the sheaf's cohomology is nonzero independent of that data), and abstains from synthesis when said obstruction class is nonzero, emitting its support as a conflict witness. \emph{[Draft language pending patent counsel review --- see \S\ref{sec:patent:defense} discussion note; not yet filed.]}
\end{patentclaim}

\begin{patentdiscussion}
Prior art (rerankers, RAG, GraphRAG) never enforces evidence closure: they always return top-$k$ for synthesis. The mandatory gate with calibrated threshold is a structural invariant, not a heuristic.

Claims 8--10 narrow the invention away from the ``combination of known components'' objection that a bare hybrid-search-plus-reranker claim invites. Claim 8 recites a specific cascade discipline --- irreversible discard, channel-distinct stages, cost ordering --- with a stated quantitative consequence, rather than generic multi-stage filtering. Claim 9 recites a construction procedure for the null distribution, which is the non-obvious element: the displacement conditions are what make the estimate valid, and no retrieval prior art constructs relevance nulls at all. Claim 10 recites a topological consistency criterion that is not reducible to pairwise contradiction checking, since cyclic inconsistency is invisible to any pairwise test.
\end{patentdiscussion}

\section{Summary}
\label{sec:retrieval:summary}

The retrieval pipeline transforms a query into a verified, mutually consistent evidence set:
\begin{enumerate}
  \item \textbf{Hybrid Search}: Vector (HNSW) + Keyword (BM25) + Graph (PPR/adjacency)
  \item \textbf{Proofreading Cascade}: $n$-stage irreversible verification over distinct evidence channels; error suppression multiplicative in depth
  \item \textbf{Calibration}: label-free FAR from an empirical null (\S\ref{sec:retrieval:far}), complemented by conformal coverage when labels exist
  \item \textbf{Consistency Gate}: sheaf cohomology over the surviving set; $H^1 \neq 0 \Rightarrow$ abstain with conflict witness
  \item \textbf{Context Selection}: deterministic top-$k$ baseline; budgeted concept coverage remains feature-gated pending answer-quality A/B
  \item \textbf{WC/1 Transport}: lossless evidence text, compact integer handles, host-side occurrence/span resolution
  \item \textbf{Forced Citations}: Synthesis only from verified units, followed by claim support and coverage checks
  \item \textbf{Evidence Closure}: cohomological form (Def.~\ref{def:closure-cohomological}); every citation traces to a verified unit admitting a consistent global reading
\end{enumerate}
This pipeline is the enforcement mechanism for LLMWiki's core invariant: \emph{answers are evidence-bounded}. The strengthening over the volume's opening statement is deliberate --- evidence closure begins as a provenance property and ends as a verifiable topological one.

%===============================================================================
% End of Volume 04
%===============================================================================
